\subsection{伴随方程}

仍然假设空间 E 是 C 上的 n 维空间，设 \(E^{*}\) 是它的对偶 (A, II, p.~40)，这是一个 C 上的 n 维空间 (Alg., II, p.~299, th. 4)；典范双线性型 \(\langle x, x^{*}\rangle\) （定义在 \(E \times E^{*}\) 上）(Alg., II, p.~234) 在这个乘积空间上连续（因为它是 \(x \in E\) 与 \(x^{*} \in E^{*}\) 的分量的多项式）。

给定齐次线性方程 (4) (IV, p.~183)，其中 \(t \mapsto A(t)\) 是 \(\mathbf{J}\) 到 \(\mathcal{L}(\mathrm{E})\) 中的一个规则映射，我们要考察是否存在一个映射 \(t \mapsto \mathbf{v}(t)\) ，它把 \(\mathbf{J}\) 映入 \(\mathrm{E}^*\) ，是 \(\mathbf{J}\) 上某个规则函数的原函数，并且使得数量函数 \(t \mapsto \langle \mathbf{u}(t), \mathbf{v}(t) \rangle\) 在 \(\mathbf{J}\) 上为常数（当 \(\mathbf{u}\) 是 (4) 的任意解时）；这等同于写出：这个函数的导数在 \(\mathbf{u}\) 与 \(\mathbf{v}\) 都可导的每一点处为零，也就是说，必须有

\[
\left\langle \frac {d {\bf u}}{d t}, {\bf v} (t) \right\rangle + \left\langle {\bf u} (t), \frac {d {\bf v}}{d t} \right\rangle = 0
\]

在这样的点处。

现在，由 (4) 有 \(\left\langle\frac{d\mathbf{u}}{dt},\mathbf{v}(t)\right\rangle=\langle A(t).\mathbf{u}(t),\mathbf{v}(t)\rangle=-\langle\mathbf{u}(t),B(t).\mathbf{v}(t)\rangle\) ，其中 \(-B(t)\) 是 \(A(t)\) 的转置 (Alg., II, p.~234)。于是 v 所必须满足的关系可以写成

\[
\left\langle \mathbf {u} (t), \frac {d \mathbf {v}}{d t} - B (t). \mathbf {v} (t) \right\rangle = 0
\]

在 \(A(t)\) 连续且 \(\mathbf{v}(t)\) 可导的所有点处。现在，对这样的点 \(t\) 以及任意的点 \(\mathbf{x}_0 \in E\) ，由 IV, p.~179 的定理 1，存在 (4) 的一个解 \(\mathbf{u}\) 使得 \(\mathbf{u}(t) = \mathbf{x}_0\) ；因此必须有 \(\left\langle \mathbf{x}_0, \frac{d\mathbf{v}}{dt} - B(t).\mathbf{v}(t) \right\rangle = 0\) （对所有的 \(\mathbf{x}_0 \in E\) ），这就蕴含 \(\frac{d\mathbf{v}}{dt} - B(t).\mathbf{v}(t) = 0\) 。因此：

命题 6. 映射 \(t \mapsto \mathbf{v}(t)\) （由 \(\mathbf{J}\) 到 \(\mathrm{E}^*\) ，且为 \(\mathbf{J}\) 上某个规则函数的原函数）使得 \(\langle \mathbf{u}(t), \mathbf{v}(t) \rangle\) 在 \(\mathbf{J}\) 上对 IV, p.~183 的方程 (4) 的每个解 \(\mathbf{u}\) 均为常数，当且仅当 \(\mathbf{v}\) 是齐次线性方程

\[
{\frac {d \mathbf {x}}{d t}} = B (t). \mathbf {x}\tag{16}
\]

的解，其中 \(-B(t)\) 是 \(A(t)\) 的转置。

方程 (16) 称为 (4) 的伴随方程；显然 (4) 也是 (16) 的伴随方程。由于矩阵 \(B(t)\) 的元素是 t 在 J 上的规则函数，前面关于线性方程得到的结果都适用于 (16)。特别地，(16) 的取值 \(x_{0}^{*}\) 于点 \(t_{0}\) 的积分可以写成 \(H(t, t_{0})\) . \(x_{0}^{*}\) ，其中 \(H(t, t_{0})\) 是 \(E^{*}\) 到其自身上的一个双射线性映射，它等同于方程

\[
\frac {d V}{d t} = B (t) V\tag{17}
\]

的取值 \(I\) 于点 \(t_0\) 的积分。由此得到（采用 IV, p.~179 的记号）

\[
\left\langle C (t, t _ {0}). \mathbf {x} _ {0}, H (t, t _ {0}). \mathbf {x} _ {0} ^ {*} \right\rangle = \left\langle \mathbf {x} _ {0}, \mathbf {x} _ {0} ^ {*} \right\rangle
\]

（对任意的 \(x_{0} \in E\) 与 \(x_{0}^{*} \in E^{*}\) ），这表明

\[
H (t, t _ {0}) = \check {C} (t, t _ {0})\tag{18}
\]

（即 \(C(t, t_{0})\) 的逆步映射）。特别地，若已知伴随方程 (16) 的一个基本积分组，则矩阵 \(H(t, t_{0})\) 就被确定，从而 \(C(t, t_{0})\) 也随之确定，于是方程 (4) 的全部积分也都随之确定。

注. 设 E 与 F 是 R（或 C）上的两个完备赋范空间，\((\mathbf{x}, \mathbf{y}) \mapsto \langle \mathbf{x}, \mathbf{y} \rangle\) 是 \(E \times F\) 上的一个连续双线性型，使得关系 `` \(\langle x, y \rangle = 0\) 对所有 \(y \in F\) ''（相应地，`` \(\langle x, y \rangle = 0\) 对所有 \(x \in E\) ''）蕴含 x = 0（相应地，y = 0）。再假设对每个 \(t \in J\) 存在一个 F 到其自身的连续线性映射 \(B(t)\) ，使得 \(\langle A(t).x, y \rangle + \langle x, B(t).y \rangle = 0\) 对每个 \((\mathbf{x}, \mathbf{y}) \in E \times F\) 成立。在这些条件下，与前面一样可以看出，为使一个把 J 映入 F 的映射 \(t \mapsto v(t)\) （某个规则函数的原函数）满足 \(\langle u(t), v(t) \rangle\) 对 (4) 的每个积分 u 都为常数，必须且只需 v 是 (16) 的一个积分，我们仍称 (16) 为 (4) 的伴随方程。

